\reading{MR-7}{Beyond Exceedance-Based Backtesting of VaR Models (Probability Integral Transform)}{grind}{4}
% ==========================================================

\begin{mrkeybox}[title={Learning objectives in this reading}]
\small
\begin{enumerate}[label=\textbf{\alph*.},leftmargin=1.7em]
\item Identify the properties of an exceedance-based backtest that indicate a VaR model is
      accurate, and describe how these properties are reflected in a PIT-based backtest.
\item Explain how to derive probability integral transforms (PITs) in the context of validating
      a VaR model.
\item Describe how the shape of the distribution of PITs can be used as an indicator of the
      quality of a VaR model.
\item Describe backtesting using PITs, and compare the various goodness-of-fit tests that can be
      used to evaluate the distribution of PITs: the Kolmogorov-Smirnov test, the
      Anderson-Darling test, and the Cram\'er-von Mises test.
\end{enumerate}
\end{mrkeybox}

% ---------------------------------------------------------
\lo{MR-7 a}{Identify the properties of an exceedance-based backtest that indicate a VaR model is accurate, and describe how these properties are reflected in a PIT-based backtest.}

\subsection{The exceedance-based approach}
Exceedance-based backtesting is a crucial statistical procedure in validating Value-at-Risk
models. It involves comparing realized outcomes against the model's forecasted values. The
primary aims are to:
\begin{itemize}
\item ensure the model's accuracy for internal risk management;
\item satisfy regulatory capital calculation requirements.
\end{itemize}

\begin{mrkeybox}[title={The three properties of a well-calibrated VaR model}]
\begin{enumerate}
\item \term{Unconditional coverage} --- ensures that the model correctly estimates the expected
      probability of exceedances.
\item \term{Independence} --- confirms that exceedances occur randomly and are not clustered.
\item \term{Conditional coverage} --- combines both unconditional coverage and independence to
      verify overall accuracy.
\end{enumerate}
\end{mrkeybox}

\subsection{1. Unconditional coverage}
Unconditional coverage checks whether the actual number of exceedances (days where losses exceed
VaR estimates) aligns with the expected number. Unlike other tests, it does not consider the
timing or independence of exceedances --- only the total count over a period.

\begin{mrtrapbox}[title={Reading the direction}]
\begin{itemize}
\item \term{Too few} exceedances $\rightarrow$ the VaR model \term{overestimates} risk
      (conservative).
\item \term{Too many} exceedances $\rightarrow$ the VaR model \term{underestimates} risk
      (aggressive).
\end{itemize}
\end{mrtrapbox}

\subsubsection{Exceedance probability in a 95\% VaR model}
\begin{enumerate}
\item A properly calibrated 95\% VaR model should have a 5\% probability of exceedance on any
      given day. This probability is 1 minus the VaR confidence level. For a 95\% VaR,
      $p = 0.05$.
\item But if you wish to find \textit{exactly} how many exceedances, then the probability of
      observing $x$ exceedances in $n$ total observations follows a binomial distribution.
\end{enumerate}

\begin{mrfmlbox}
\[ p(x) \;=\; \frac{n!}{(n-x)!\;x!}\,(1-p)^{\,n-x}\, p^{\,x} \]
\mrvk{$n$ = total number of observations;\quad $x$ = number of exceedances;\quad $p$ = the
per-day exceedance probability, i.e.\ $1 -$ confidence level.}{}
\end{mrfmlbox}

\begin{mrexbox}[title={Example --- exactly three exceedances}]
For 100 observations with a 5\% exceedance probability, the probability of exactly 3
exceedances is:
\tcblower
\small
\[ p(3) \;=\; \frac{100!}{(100-3)!\;3!}\,(1-0.05)^{97}\,(0.05)^{3} \;=\; \mathbf{13.96\%} \]
There is a 13.96\% chance of observing exactly 3 exceedances in 100 days.
\end{mrexbox}

\begin{center}
\begin{tikzpicture}
\begin{axis}[
  width=11.4cm, height=5.4cm,
  xmin=-0.8, xmax=14.8, ymin=0, ymax=0.21,
  xlabel={\scriptsize Number of exceedances in 100 days},
  ylabel={\scriptsize Likelihood},
  xlabel style={font=\scriptsize}, ylabel style={font=\scriptsize},
  tick label style={font=\scriptsize},
  xtick={0,1,2,3,4,5,6,7,8,9,10,11,12,13,14},
  ytick={0,0.05,0.10,0.15,0.20},
  yticklabel style={/pgf/number format/fixed},
  axis line style={draw=black!55}, grid=none, clip=false]
\addplot[ybar, bar width=8pt, fill=PaleNavy, draw=FRMNavy!60, line width=0.3pt] coordinates {
(0,0.0059)(1,0.0312)(2,0.0812)(4,0.1781)(5,0.1800)(6,0.1500)(7,0.1060)(8,0.0649)
(9,0.0349)(10,0.0167)(11,0.0072)(12,0.0028)(13,0.0010)(14,0.0003)};
\addplot[ybar, bar width=8pt, fill=FRMGold!45, draw=FRMGold, line width=0.5pt] coordinates {(3,0.1396)};
\draw[FRMRed, dashed, line width=0.7pt] (axis cs:5,0) -- (axis cs:5,0.192);
\node[anchor=west, font=\scriptsize\sffamily, text=FRMRed, fill=white, inner sep=1.5pt] at (axis cs:5.4,0.186)
  {expected $= np = 5$};
\node[anchor=west, font=\scriptsize\sffamily, text=FRMGold, fill=white, inner sep=1.5pt] at (axis cs:0.2,0.155)
  {$p(3) = 13.96\%$};
\draw[FRMGold, -{Stealth[length=3.5pt]}, line width=0.5pt]
  (axis cs:2.3,0.152) -- (axis cs:2.85,0.142);
\end{axis}
\end{tikzpicture}
\end{center}
\figcap{The whole binomial distribution of exceedance counts, not just the one number the
example asks for. The mode sits at 4 and 5, so counts anywhere in the 2 to 8 range are
unremarkable; it is the thin shoulders beyond that which signal a model problem.}

\subsubsection{VaR model validation using Kupiec unconditional coverage}
\begin{itemize}
\item The goal is to validate whether a VaR model correctly estimates risk by analyzing the
      number of exceedances (days when losses exceed the VaR estimate).
\item Hypothesis testing is used to check if the actual exceedances align with expectations.
\end{itemize}

The test follows a chi-squared distribution with one degree of freedom. At a 95\% confidence
level, the critical value is \term{3.84}. If the LR test statistic exceeds this value, the VaR
model is rejected as inaccurate.

\begin{mrfmlbox}
\[ LR_{uc} \;=\; -2\ln\!\left[ (1-p)^{\,n-x} p^{\,x} \right]
   \;+\; 2\ln\!\left\{ \left[ 1 - (x/n) \right]^{\,n-x} (x/n)^{\,x} \right\} \]
\mrvk{$p$ = the exceedance probability the model claims;\quad $x/n$ = the exceedance rate actually
observed. The first term is the likelihood under the model, the second under the data.}{}
\end{mrfmlbox}

\subsection{2. Independence}
We need to check if VaR exceedances are independent or if they show \term{clustering} --- that
is, if an exceedance today increases the probability of another exceedance tomorrow. If
clustering occurs, it means the model is underestimating risk because it fails to adjust for
changing market conditions. An accurate VaR model should ensure that exceedances occur with the
expected probability (e.g.\ 5\% for a 95\% VaR).

\subsubsection{Using Markov chains to analyze independence}
Exceedances can be modeled as a \term{Markov chain}, where future states depend on past states.

\begin{center}
\begin{tikzpicture}[every node/.style={font=\scriptsize\sffamily}]
  \node[circle, draw=FRMNavy, fill=PaleNavy, line width=0.9pt, minimum size=2.1cm,
        align=center, text=black] (n0) at (0,0) {\textbf{State 0}\\ No\\ exceedance};
  \node[circle, draw=FRMRed, fill=PaleRed, line width=0.9pt, minimum size=2.1cm,
        align=center, text=black] (n1) at (6.4,0) {\textbf{State 1}\\ Exceedance};

  \draw[FRMNavy, -{Stealth[length=5pt]}, line width=0.8pt]
    (n0) to[bend left=22] node[above, text=FRMNavy] {$\pi_{01}$} (n1);
  \draw[FRMRed, -{Stealth[length=5pt]}, line width=0.8pt]
    (n1) to[bend left=22] node[below, text=FRMRed] {$\pi_{10}$} (n0);
  \draw[FRMNavy, -{Stealth[length=5pt]}, line width=0.8pt]
    (n0) to[out=150, in=210, looseness=5]
    node[left, text=FRMNavy] {$\pi_{00}$} (n0);
  \draw[FRMRed, -{Stealth[length=5pt]}, line width=0.8pt]
    (n1) to[out=30, in=-30, looseness=5]
    node[right, text=FRMRed] {$\pi_{11}$} (n1);
\end{tikzpicture}
\end{center}
\figcap{$\pi_{01}$ is the chance of an exceedance tomorrow given none today; $\pi_{11}$ is the
chance of an exceedance tomorrow given one today. If those two differ, yesterday tells you
something about today, and the exceedances are not independent.}

\begin{mrdefbox}[title={The independence hypothesis}]
\textbf{$H_0$: $\pi_{01} = \pi_{11}$} --- exceedances are independent. The probability of an
exceedance today is the same whether or not there was an exceedance yesterday.

\medskip
Note that the test statistic for independence, $LR_{ind}$, also follows a chi-squared
distribution with one degree of freedom.
\end{mrdefbox}

\subsection{3. Conditional coverage}
Conditional coverage combines unconditional coverage (whether the number of exceedances matches
expectations) and independence (whether exceedances occur randomly or show clustering). It
checks two things simultaneously:
\begin{itemize}
\item Does the model correctly predict the expected exceedance rate?
\item Are exceedances independent over time?
\end{itemize}

\begin{mrfmlbox}
\[ LR_{cc} \;=\; LR_{uc} \;+\; LR_{ind} \]
\mrvk{each component is $\chi^2$ with one degree of freedom, so the sum is $\chi^2$ with two.}{}
\end{mrfmlbox}

\begin{mrtrapbox}[title={Two critical values, two different tests}]
Using a 95\% confidence level, rejection of the \textbf{model} occurs if
$LR_{cc} > \mathbf{5.99}$, and rejection of the \textbf{independence term alone} occurs if
$LR_{ind} > \mathbf{3.84}$.
\end{mrtrapbox}

% ---------------------------------------------------------
\lo{MR-7 b}{Explain how to derive probability integral transforms (PITs) in the context of validating a VaR model.}

\begin{center}
\small
\begin{tabularx}{\linewidth}{@{}p{3.6cm} >{\raggedright\arraybackslash}X@{}}
\arrayrulecolor{FRMRule}\toprule
\rowcolor{FRMNavy} \hdr{Approach} & \hdr{What it looks at}\\
\textbf{Exceedance-based} & Only focuses on whether losses exceed a preselected VaR threshold
  (e.g.\ the 95\% or 99\% confidence level). It does \textit{not} assess how well the entire
  forecasted distribution aligns with actual returns.\\
\rowcolor{FRMSoft}
\textbf{PIT} & Instead of just counting exceedances, the PIT approach evaluates how well the
  predicted distribution fits the realized outcomes. It transforms actual profit and loss
  observations onto a standardized uniform scale to assess the model's accuracy.\\
\bottomrule
\end{tabularx}
\end{center}

VaR estimation predicts potential outcome distributions, focusing on tail risks. In a static
distribution, straightforward statistical analysis would suffice. However, financial portfolios
experience \term{dynamic} distributions. The PIT approach standardizes P\&L observations across
periods, allowing unbiased validation.

\subsection{The four steps, with a worked observation}

\begin{center}
\begin{tikzpicture}[every node/.style={font=\scriptsize\sffamily},
  stepbox/.style={rectangle, rounded corners=2.5pt, draw=FRMNavy!55, fill=PaleNavy,
    line width=0.7pt, text width=2.85cm, align=center, inner sep=5pt, minimum height=1.5cm}]
  \node[stepbox] (s1) at (0,0)
    {\textbf{Step 1}\\ Forecast the distribution for period $T$};
  \node[stepbox, right=0.55cm of s1] (s2)
    {\textbf{Step 2}\\ Collect the actual observation for period $T$};
  \node[stepbox, right=0.55cm of s2] (s3)
    {\textbf{Step 3}\\ Read where it falls on the forecast CDF};
  \node[stepbox, right=0.55cm of s3] (s4)
    {\textbf{Step 4}\\ Analyse the PITs for uniformity and independence};
  \foreach \a/\b in {s1/s2, s2/s3, s3/s4}{
    \draw[FRMGold, -{Stealth[length=4.5pt]}, line width=1pt] (\a) -- (\b);}
  \node[below=2pt of s1, text width=2.85cm, align=center, text=black!65]
    {P\&L $\sim N(\$0,\ \$500)$};
  \node[below=2pt of s2, text width=2.85cm, align=center, text=black!65]
    {Actual P\&L $= -\$250$};
  \node[below=2pt of s3, text width=2.85cm, align=center, text=black!65]
    {$-0.5\sigma \Rightarrow$ \textbf{PIT} $=0.308$};
  \node[below=2pt of s4, text width=2.85cm, align=center, text=black!65]
    {Same three properties as before};
\end{tikzpicture}
\end{center}

\begin{center}
\begin{tikzpicture}
\begin{axis}[name=pdfax, width=6.6cm, height=4.2cm,
  xmin=-3.2, xmax=3.2, ymin=0, ymax=0.46,
  axis y line=none, axis x line=bottom, axis line style={draw=black!55},
  xtick={-2,-0.5,1,2},
  xticklabels={$-1000$,$-250$,$500$,$1000$},
  tick label style={font=\tiny}, clip=false,
  title={\scriptsize\sffamily\bfseries\color{FRMNavy}Forecast P\&L distribution}]
\addplot[domain=-3.2:3.2, samples=200, smooth, black!80, line width=0.9pt] {0.39894*exp(-x^2/2)};
\addplot[domain=-3.2:-0.5, samples=120, smooth, draw=none, fill=FRMGold!40]
  {0.39894*exp(-x^2/2)} \closedcycle;
\draw[FRMGold, line width=0.9pt] (axis cs:-0.5,0) -- (axis cs:-0.5,0.352);
\node[anchor=west, font=\tiny\sffamily, text=FRMGold] at (axis cs:-3.1,0.40)
  {area $=30.8\%$};
\end{axis}
\begin{axis}[name=cdfax, at={($(pdfax.east)+(1.5cm,0)$)}, anchor=west,
  width=6.6cm, height=4.2cm,
  xmin=-3.2, xmax=3.2, ymin=0, ymax=1.05,
  axis y line=left, axis x line=bottom, axis line style={draw=black!55},
  xtick={-2,-0.5,1,2}, xticklabels={$-1000$,$-250$,$500$,$1000$},
  ytick={0,0.308,0.5,1}, yticklabels={0,0.308,0.5,1},
  tick label style={font=\tiny}, clip=false,
  title={\scriptsize\sffamily\bfseries\color{FRMNavy}Same forecast, as a CDF}]
\addplot[domain=-3.2:3.2, samples=200, smooth, black!80, line width=0.9pt]
  {1/(1+exp(-1.702*x))};
\draw[FRMGold, dashed, line width=0.8pt] (axis cs:-0.5,0) -- (axis cs:-0.5,0.308);
\draw[FRMGold, dashed, line width=0.8pt] (axis cs:-3.2,0.308) -- (axis cs:-0.5,0.308);
\node[circle, fill=FRMGold, inner sep=1.6pt] at (axis cs:-0.5,0.308) {};
\node[anchor=west, font=\tiny\sffamily, text=FRMGold] at (axis cs:-0.35,0.28)
  {PIT $=0.308$};
\end{axis}
\end{tikzpicture}
\end{center}
\figcap{The PIT is just the observation read off the vertical axis of the forecast CDF. The
shaded area on the left is the same 30.8\%, seen as a probability instead of a height. A model
that is right on average will scatter these readings evenly over $[0,1]$.}

\begin{mrkeybox}[title={PIT-based backtesting and its properties}]
\begin{itemize}
\item \term{Uniformity:} PITs should be uniformly distributed over the interval $[0,1]$ if the
      risk is accurately modeled. This reflects the \textit{unconditional coverage} property,
      where each P\&L outcome corresponds to a uniformly distributed probability.
\item \term{Independence:} the PIT series should also be independently and identically
      distributed (i.i.d.), validating the model's \textit{independence} property.
\end{itemize}
\end{mrkeybox}

% ---------------------------------------------------------
\lo{MR-7 c}{Describe how the shape of the distribution of PITs can be used as an indicator of the quality of a VaR model.}

PITs offer valuable insights into the calibration of a VaR model. By plotting PITs for all
observations, a well-calibrated model should exhibit a uniform distribution over $[0,1]$. This
uniformity signifies that the model's predicted probabilities align closely with the observed
outcomes, ensuring consistency across the entire spectrum of risk scenarios.

\begin{itemize}
\item If the VaR model perfectly forecasts the next period's P\&L distribution, the resulting
      PITs will be uniformly distributed between 0 and 1.
\item This means each decile (0--0.1, 0.1--0.2, and so on) will contain 10\% of the PITs.
\end{itemize}

\begin{center}
\begin{tikzpicture}
\begin{axis}[name=perf, width=7.0cm, height=4.4cm,
  xmin=0.3, xmax=10.7, ymin=0, ymax=0.155,
  xlabel={\scriptsize Decile band}, ylabel={\scriptsize Portion of observations},
  xlabel style={font=\scriptsize}, ylabel style={font=\scriptsize},
  tick label style={font=\tiny},
  xtick={1,2,3,4,5,6,7,8,9,10},
  xticklabels={0,.1,.2,.3,.4,.5,.6,.7,.8,.9},
  ytick={0,0.05,0.10,0.15}, yticklabels={0\%,5\%,10\%,15\%},
  axis line style={draw=black!55},
  title={\scriptsize\sffamily\bfseries\color{FRMGreen}Perfect forecasting}]
\addplot[ybar, bar width=11pt, fill=PaleGreen, draw=FRMGreen!70, line width=0.3pt] coordinates {
(1,0.10)(2,0.10)(3,0.10)(4,0.10)(5,0.10)(6,0.10)(7,0.10)(8,0.10)(9,0.10)(10,0.10)};
\addplot[FRMGreen, dashed, line width=0.8pt, forget plot, sharp plot]
  coordinates {(0.3,0.10)(10.7,0.10)};
\end{axis}
\begin{axis}[name=imp, at={($(perf.east)+(1.9cm,0)$)}, anchor=west,
  width=7.0cm, height=4.4cm,
  xmin=0.3, xmax=10.7, ymin=0, ymax=0.155,
  xlabel={\scriptsize Decile band},
  xlabel style={font=\scriptsize},
  tick label style={font=\tiny},
  xtick={1,2,3,4,5,6,7,8,9,10},
  xticklabels={0,.1,.2,.3,.4,.5,.6,.7,.8,.9},
  ytick={0,0.05,0.10,0.15}, yticklabels={0\%,5\%,10\%,15\%},
  axis line style={draw=black!55},
  title={\scriptsize\sffamily\bfseries\color{FRMRed}Imperfect forecasting}]
\addplot[ybar, bar width=11pt, fill=PaleRed, draw=FRMRed!70, line width=0.3pt] coordinates {
(1,0.13)(2,0.09)(3,0.10)(4,0.09)(5,0.10)(6,0.10)(7,0.11)(8,0.07)(9,0.08)(10,0.13)};
\addplot[FRMRed, dashed, line width=0.8pt, forget plot, sharp plot]
  coordinates {(0.3,0.10)(10.7,0.10)};
\end{axis}
\end{tikzpicture}
\end{center}
\figcap{Left, every decile holds exactly 10\% of the PITs. Right, the two end bands are
overfull at 13\% while the 0.7--0.9 region is underfull. This chart illustrates a set of PITs
that indicate that forecasted distributions underestimate losses or gains in some places and
overestimate losses or gains in other places.}

\begin{center}
\begin{tikzpicture}
% ===== LEFT: the forecast versus what actually happened =====
\begin{axis}[name=fc, width=6.0cm, height=4.2cm,
  xmin=-4.2, xmax=4.2, ymin=0, ymax=0.50,
  axis y line=none, axis x line=bottom, axis line style={draw=black!55},
  xlabel={\scriptsize P\&L}, xlabel style={font=\tiny},
  xtick={\empty}, clip=false,
  title={\scriptsize\sffamily\bfseries\color{FRMNavy}Forecast is too \textit{narrow}},
  legend style={font=\tiny\sffamily, draw=FRMRule, fill=white,
    at={(0.5,-0.14)}, anchor=north, legend columns=-1, column sep=5pt}]
\addplot[FRMNavy, line width=1.2pt, smooth, domain=-4.2:4.2, samples=150]
  {0.455*exp(-(x^2)/1.0)};
\addlegendentry{model forecast}
\addplot[FRMRed, line width=1.2pt, dashed, smooth, domain=-4.2:4.2, samples=150]
  {0.30*exp(-(x^2)/3.4)};
\addlegendentry{what actually happens}
\node[anchor=south, font=\tiny\sffamily, text=FRMRed, fill=white, inner sep=1pt]
  at (axis cs:-3.0,0.055) {real outcomes};
\node[anchor=south, font=\tiny\sffamily, text=FRMRed, fill=white, inner sep=1pt]
  at (axis cs:3.0,0.055) {land out here};
\end{axis}
\draw[FRMGold, -{Stealth[length=5.5pt]}, line width=1.2pt]
  ($(fc.east)+(0.6cm,0)$) -- ($(fc.east)+(1.9cm,0)$);
\node[anchor=south, font=\tiny\sffamily, text=FRMGold, align=center]
  at ($(fc.east)+(1.25cm,0.14cm)$) {so the PITs};
% ===== RIGHT: the PIT histogram that results =====
\begin{axis}[name=ph, at={($(fc.east)+(3.0cm,0)$)}, anchor=west,
  width=6.2cm, height=4.2cm,
  xmin=0.3, xmax=10.7, ymin=0, ymax=0.165,
  xlabel={\scriptsize Decile band}, ylabel={\scriptsize Portion},
  xlabel style={font=\tiny}, ylabel style={font=\tiny},
  tick label style={font=\tiny},
  xtick={1,3,5,7,9}, xticklabels={0,.2,.4,.6,.8},
  ytick={0,0.05,0.10,0.15}, yticklabels={0\%,5\%,10\%,15\%},
  axis line style={draw=black!55}, clip=false,
  title={\scriptsize\sffamily\bfseries\color{FRMRed}pile up at \textit{both ends}}]
\addplot[ybar, bar width=9pt, fill=PaleRed, draw=FRMRed!70, line width=0.3pt] coordinates {
(1,0.13)(2,0.09)(3,0.10)(4,0.09)(5,0.10)(6,0.10)(7,0.11)(8,0.07)(9,0.08)(10,0.13)};
\addplot[FRMRed, dashed, line width=0.8pt, forget plot, sharp plot]
  coordinates {(0.3,0.10)(10.7,0.10)};
\end{axis}
\end{tikzpicture}
\end{center}
\figcap{Why the end bands are overfull. If the model's forecast distribution is narrower than
reality, then outcomes that the model treats as extreme happen more often than it expects, and
those outcomes map to PIT values very close to 0 and very close to 1. The histogram on the right
\textit{is} the mismatch on the left, counted. A forecast that is too \textbf{wide} produces the
opposite --- PITs bunching in the middle deciles and thin at the ends.}

% ---------------------------------------------------------
\lo{MR-7 d}{Describe backtesting using PITs, and compare the various goodness-of-fit tests that can be used to evaluate the distribution of PITs: the Kolmogorov-Smirnov test, the Anderson-Darling test, and the Cram\'er-von Mises test.}

\subsection{Evaluating PIT distributions}
To assess the quality of a VaR model using the shape of the PIT distribution:
\begin{itemize}
\item Use visual tools like histograms or QQ plots to detect uniformity deviations.
\item Apply goodness-of-fit tests such as Kolmogorov-Smirnov (KS), Anderson-Darling (AD), or
      Cram\'er-von Mises (CvM).
\end{itemize}

These tests assess how well the empirical distribution of PITs matches the theoretical uniform
distribution. They often utilize QQ plots, which visually compare the quantiles of the two
distributions.

\begin{mrtrapbox}[title={Remember}]
Deviations from the straight diagonal line in a QQ plot suggest a mismatch.
\end{mrtrapbox}

\subsection{1) Kolmogorov-Smirnov test}
\begin{itemize}
\item The KS test is a widely used non-parametric test that evaluates goodness-of-fit by
      comparing the empirical distribution function (EDF) of a sample with the reference uniform
      distribution function over $[0,1]$.
\item It measures the \term{maximum distance} between the EDF and the reference CDF. This
      maximum difference is the KS statistic, $D$.
\item The KS test is ideal for checking the overall uniformity of PIT distributions but may miss
      tail-specific misspecifications. It is commonly used for quick checks of distributional
      uniformity.
\end{itemize}

\begin{mrfmlbox}
\[ D \;=\; \max_j \left| z_j - A_j \right| \]
\mrvk{$j$ = ranked observation out of $n$ total observations;\quad $z_j$ = theoretical (perfect)
distribution;\quad $A_j$ = actual CDF distribution.}{%
\par\smallskip\textbf{\sffamily Remember:} a \textit{smaller} KS statistic $D$ indicates a
better fit.}
\end{mrfmlbox}

\begin{center}
\begin{tikzpicture}
\begin{axis}[
  width=8.2cm, height=5.6cm,
  xmin=0, xmax=1, ymin=0, ymax=1,
  xlabel={\scriptsize PIT value}, ylabel={\scriptsize Portion of observations $\le$ PIT},
  xlabel style={font=\scriptsize}, ylabel style={font=\scriptsize},
  tick label style={font=\scriptsize},
  xtick={0,0.25,0.5,0.75,1}, ytick={0,0.25,0.5,0.75,1},
  axis line style={draw=black!55}, clip=false]
\addplot[FRMNavy, line width=1.2pt, domain=0:1, samples=2] {x};
\addplot[black!70, line width=0.9pt, smooth, domain=0:1, samples=120]
  {x + 0.19*sin(deg(3.14159*x))*cos(deg(0.9*3.14159*x))};
\draw[FRMRed, {Stealth[length=4pt]}-{Stealth[length=4pt]}, line width=0.9pt]
  (axis cs:0.245,0.245) -- (axis cs:0.245,0.408);
\node[anchor=west, font=\scriptsize\sffamily, text=FRMRed] at (axis cs:0.27,0.33) {$D$};
\node[anchor=north west, font=\scriptsize\sffamily, text=FRMNavy, fill=white, inner sep=1.5pt]
  at (axis cs:0.62,0.44) {perfect uniform};
\node[anchor=south east, font=\scriptsize\sffamily, text=black!70, fill=white, inner sep=1.5pt]
  at (axis cs:0.42,0.62) {actual EDF};
\end{axis}
\end{tikzpicture}
\end{center}
\figcap{$D$ is the single widest vertical gap anywhere between the two curves. Because it takes
only the maximum, a model can be badly wrong in the far tail and still score well if the gap
there is narrower than the gap in the middle. That is the KS test's blind spot.}

\subsection{2) Anderson-Darling test}
Similar to the KS test, but places \textit{more weight} on discrepancies in the \textit{tails}
of the distribution. This makes it more suitable for risk management applications. The AD test
places higher emphasis on tail observations, making it a more appealing option for risk
managers.

\begin{mrfmlbox}
\[ A^{2} \;=\; -n \;-\; \frac{1}{n}\sum_{j=1}^{n} (2j-1)\Bigl[ \ln(z_j) +
   \ln(1 - z_{\,n+1-j}) \Bigr] \]
\mrvk{$n$ = number of observations;\quad $j$ = rank of the observation;\quad $z_j$ = the $j$th
ordered PIT value. The $\ln$ terms blow up as $z_j$ approaches 0 or 1, which is precisely why
tail discrepancies are weighted so heavily.}{}
\end{mrfmlbox}

\subsection{3) Cram\'er-von Mises test}
The CvM test is very similar to the KS test, but instead of evaluating the \textit{maximum}
difference between distributions, it considers the \textit{mean squared deviation}.

\begin{mrfmlbox}
\[ W^{2} \;=\; \sum_{j=1}^{n} \left[ z_j - \frac{(2j-1)}{2n} \right]^{2} + \frac{1}{12n} \]
\mrvk{$z_j$ = the $j$th ordered PIT value;\quad $(2j-1)/2n$ = where that observation would sit
under a perfect uniform;\quad $n$ = number of observations.}{}
\end{mrfmlbox}

\subsection{The three tests side by side}

\begin{center}
\small
\begin{tabularx}{\linewidth}{@{}p{2.1cm} p{2.9cm} >{\raggedright\arraybackslash}X >{\raggedright\arraybackslash}X@{}}
\arrayrulecolor{FRMRule}\toprule
\rowcolor{FRMNavy} \hdr{Test} & \hdr{What it measures} & \hdr{Strengths} & \hdr{Limitations}\\
\midrule
\textbf{Kolmogorov-\newline Smirnov (KS)}
  & The \textit{maximum} distance between the EDF and the reference CDF, $D$.
  & Simple and widely used for general goodness-of-fit evaluation. Detects deviations in the
    \textit{central} portion of the distribution effectively.
  & Less sensitive to deviations in the \textit{tails} of the distribution.\\
\rowcolor{FRMSoft}
\textbf{Anderson-\newline Darling (AD)}
  & The same comparison, but weighted so that tail observations count for more.
  & Corrects the KS test's lack of sensitivity to tail deviations. Focuses on the tails, giving
    more statistical power for extreme quantiles (e.g.\ 99\% confidence levels).
  & More computationally intensive compared to the KS test.\\
\textbf{Cram\'er-von\newline Mises (CvM)}
  & The \textit{mean squared deviation} across the whole distribution rather than the single
    largest gap.
  & Uses every observation rather than only the worst one, so a pattern of small
    misfits still registers.
  & Like KS, it does not single out the tails.\\
\bottomrule
\end{tabularx}
\end{center}
\figcap{All three compare the same two curves. They differ only in how they summarise the gap
between them: KS takes the largest gap, CvM averages the squared gaps, and AD averages them
with the tails weighted up.}

\subsection{Independence of PIT values}
Berkowitz (2001) describes an approach for how to evaluate independence through PIT values. He
notes that it may be easier to conduct independence tests \textit{outside} of the range of a
uniform distribution. For this test, evaluating the independence of PIT values requires
transforming PITs into \term{standard normal variables}.

\begin{center}
\small
\begin{tabular}{c r r r r}
\arrayrulecolor{FRMRule}\toprule
\rowcolor{FRMNavy} \hdr{Day} & \hdr{Expected std.\ deviation} & \hdr{Actual P\&L} & \hdr{PIT}
  & \hdr{Standard normal equivalent}\\
1 & 2.00\% & 4.24\% & 0.983 & 2.12\\
\rowcolor{FRMSoft} 2 & 2.25\% & 0.95\% & 0.664 & 0.42\\
3 & 1.65\% & $-0.70\%$ & 0.335 & $-0.43$\\
\rowcolor{FRMSoft} 4 & 1.00\% & 0.40\% & 0.656 & 0.40\\
5 & 3.25\% & $-5.48\%$ & 0.046 & $-1.68$\\
\rowcolor{FRMSoft} 6 & 2.50\% & $-0.65\%$ & 0.397 & $-0.26$\\
7 & 1.85\% & $-3.03\%$ & 0.051 & $-1.64$\\
\bottomrule
\end{tabular}
\end{center}
\figcap{The same P\&L is expressed three ways across the row: as a raw return, as a PIT on
$[0,1]$, and finally as a $z$-score. Day 5's $-5.48\%$ against an expected 3.25\% standard
deviation lands at a PIT of 0.046, deep in the left tail.}

The series in the last column can then be analyzed for independence across observations using
well-known log-likelihood functions.

Lastly, \term{conditional coverage} can be visually checked by viewing scatterplots of PITs
versus lagged PITs. Ideally, there should be no clustering present, and a scatterplot should
reveal a uniform distribution across all subportfolios.

\begin{center}
\begin{tikzpicture}
\begin{axis}[
  width=6.6cm, height=6.0cm,
  xmin=0, xmax=1, ymin=0, ymax=1,
  xlabel={\scriptsize Lagged PIT}, ylabel={\scriptsize PIT},
  xlabel style={font=\scriptsize}, ylabel style={font=\scriptsize},
  tick label style={font=\scriptsize},
  xtick={0,0.25,0.5,0.75,1}, ytick={0,0.25,0.5,0.75,1},
  axis line style={draw=black!55}]
\addplot[only marks, mark=*, mark size=0.55pt, color=FRMNavy!65] coordinates {
(0.238,0.544)(0.37,0.604)(0.626,0.066)(0.013,0.837)(0.259,0.234)(0.996,0.47)(0.836,0.476)
(0.639,0.151)(0.635,0.868)(0.523,0.741)(0.671,0.064)(0.758,0.591)(0.301,0.031)(0.866,0.473)
(0.719,0.879)(0.714,0.921)(0.395,0.801)(0.445,0.936)(0.879,0.097)(0.136,0.217)(0.965,0.436)
(0.627,0.301)(0.507,0.386)(0.351,0.585)(0.584,0.904)(0.682,0.929)(0.856,0.991)(0.671,0.163)
(0.861,0.965)(0.905,0.569)(0.714,0.211)(0.832,0.574)(0.285,0.063)(0.854,0.99)(0.089,0.801)
(0.41,0.151)(0.294,0.769)(0.873,0.044)(0.615,0.045)(0.718,0.331)(0.881,0.981)(0.505,0.999)
(0.31,0.077)(0.6,0.031)(0.197,0.408)(0.61,0.156)(0.042,0.868)(0.314,0.959)(0.897,0.378)
(0.46,0.52)(0.644,0.596)(0.559,0.62)(0.941,0.507)(0.431,0.72)(0.238,0.301)(0.978,0.521)
(0.548,0.011)(0.415,0.58)(0.02,0.616)(0.632,0.06)(0.627,0.466)(0.679,0.353)(0.707,0.738)
(0.022,0.061)(0.676,0.963)(0.251,0.456)(0.593,0.32)(0.364,0.313)(0.369,0.596)(0.3,0.377)
(0.772,0.027)(0.569,0.735)(0.31,0.223)(0.804,0.239)(0.187,0.435)(0.698,0.102)(0.322,0.334)
(0.834,0.438)(0.856,0.169)(0.337,0.65)(0.885,0.451)(0.225,0.121)(0.53,0.191)(0.807,0.838)
(0.184,0.279)(0.807,0.642)(0.806,0.345)(0.13,0.292)(0.794,0.271)(0.346,0.417)(0.42,0.41)
(0.921,0.156)(0.005,0.943)(0.88,0.987)(0.434,0.95)(0.927,0.222)(0.746,0.837)(0.663,0.519)
(0.289,0.341)(0.227,0.068)(0.589,0.287)(0.81,0.045)(0.904,0.694)(0.924,0.897)(0.9,0.577)
(0.013,0.745)(0.172,0.3)(0.663,0.525)(0.414,0.939)(0.612,0.341)(0.252,0.862)(0.477,0.782)
(0.352,0.197)(0.535,0.817)(0.171,0.792)(0.922,0.806)(0.823,0.008)(0.629,0.863)(0.05,0.271)
(0.269,0.527)(0.423,0.473)(0.776,0.002)(0.055,0.127)(0.125,0.068)(0.975,0.854)(0.086,0.502)
(0.316,0.315)(0.351,0.647)(0.587,0.361)(0.191,0.329)(0.124,0.556)(0.716,0.38)(0.08,0.179)
(0.373,0.604)(0.783,0.38)(0.801,0.623)(0.432,0.372)(0.496,0.703)(0.421,0.694)(0.461,0.245)
(0.536,0.695)(0.072,0.425)(0.426,0.88)(0.936,0.374)(0.898,0.791)(0.262,0.464)(0.123,0.813)
(0.662,0.887)(0.792,0.668)(0.734,0.564)(0.103,0.588)(0.005,0.144)(0.774,0.044)(0.092,0.099)
(0.88,0.179)(0.023,0.842)(0.121,0.844)(0.674,0.836)(0.952,0.579)(0.799,0.036)(0.767,0.511)
(0.715,0.107)(0.749,0.935)(0.061,0.324)(0.564,0.828)(0.242,0.18)(0.25,0.616)(0.754,0.394)
(0.367,0.397)(0.35,0.418)(0.083,0.5)(0.973,0.413)(0.747,0.161)(0.691,0.756)(0.674,0.517)
(0.484,0.643)(0.897,0.149)(0.096,0.748)(0.917,0.517)(0.443,0.719)(0.186,0.267)(0.199,0.586)
(0.315,0.232)(0.691,0.953)(0.296,0.705)(0.413,0.854)(0.585,0.267)(0.218,0.023)(0.479,0.383)
(0.172,0.36)(0.322,0.774)(0.144,0.991)(0.48,0.599)(0.468,0.835)(0.822,0.557)(0.481,0.721)
(0.857,0.4)(0.734,0.96)(0.467,0.23)(0.235,0.718)(0.675,0.959)(0.854,0.242)(0.19,0.259)
(0.187,0.705)(0.859,0.9)(0.255,0.865)(0.313,0.423)(0.729,0.086)(0.093,0.834)(0.292,0.357)
(0.58,0.676)(0.007,0.335)(0.436,0.486)(0.21,0.585)(0.955,0.391)(0.544,0.119)(0.275,0.665)
(0.113,0.887)(0.909,0.097)(0.941,0.374)(0.772,0.757)(0.296,0.676)(0.654,0.806)(0.266,0.754)
(0.961,0.673)(0.536,0.113)(0.494,0.352)(0.718,0.679)(0.566,0.182)(0.646,0.631)(0.179,0.89)
(0.655,0.123)(0.932,0.141)(0.332,0.72)(0.597,0.555)(0.647,0.458)(0.312,0.176)(0.069,0.716)
(0.754,0.543)(0.74,0.359)(0.266,0.383)(0.873,0.042)(0.505,0.247)(0.769,0.354)(0.333,0.403)
(0.541,0.772)(0.353,0.847)(0.112,0.27)(0.1,0.113)(0.779,0.727)(0.185,0.189)(0.417,0.743)
(0.816,0.749)(0.592,0.146)(0.398,0.194)(0.528,0.568)(0.202,0.25)(0.782,0.03)(0.803,0.891)
(0.949,0.383)};
\end{axis}
\end{tikzpicture}
\end{center}
\figcap{What a clean result looks like: an even cloud with no diagonal streak and no dense
corner. A band running along the diagonal would mean a high PIT tends to follow a high PIT,
which is clustering.}


% ==========================================================
